\documentclass[11pt]{article}
\usepackage[margin=1in]{geometry}
\usepackage{amsmath,amssymb}
\newcommand{\finalanswer}[1]{\par\medskip\noindent\fbox{\parbox{0.94\linewidth}{\textbf{Answer:} #1}}}
\title{STAT 412 Quiz 5 --- Question 6}
\author{}\date{}
\begin{document}\maketitle
\section*{Problem}
Rivet breaking strength has mean $6500$ psi and standard deviation $200$ psi.
For sample means, answer (a) for $n=64$, and (b) find $n$ so that
$P(6470<\bar X<6530)=0.90$.

\section*{Work}
(a) With $n=64$,
\[
\sigma_{\bar X}=\frac{200}{\sqrt{64}}=25.
\]
The endpoints are
\[
z_1=\frac{6440-6500}{25}=-2.40,\qquad
z_2=\frac{6560-6500}{25}=2.40.
\]
So
\[
P(6440<\bar X<6560)=0.9918-0.0082=0.9836.
\]
(b) The interval $6470$ to $6530$ is centered at the population mean $6500$ and
has half-width $30$:
\[
6470=6500-30,\qquad 6530=6500+30.
\]
For central probability $0.90$, the two tails total $0.10$, so each tail is
$0.05$. The upper cumulative area is $0.95$, whose normal cutoff is
$z\approx1.645$.

We want the endpoint $6530$ to be $1.645$ standard errors above the mean:
\[
30=1.645\left(\frac{200}{\sqrt n}\right).
\]
Thus
\[
\sqrt n=\frac{1.645(200)}{30},\qquad
n=\left(\frac{1.645(200)}{30}\right)^2\approx120.27.
\]
Round up, so $n=121$.

\finalanswer{(a) $0.9836$ \quad (b) $n=121$}
\end{document}
